% The same 4x4 maze twice: as a drawing of walls, and as the graph this report
% works on. The outer boundary is closed all the way round and the exit is a
% marked cell inside it, which is how the subject's own example is drawn.
\begin{figure}[htbp]
  \centering
  \begin{subfigure}[b]{0.46\textwidth}
    \centering
    \begin{tikzpicture}[scale=0.95]
      \fill[cellgrey] (0,0) rectangle (4,-4);
      \fill[exitorange!30] (2,-1) rectangle (3,-2);   % the exit cell
      \draw[wall segment] (0,0) rectangle (4,-4);     % closed on every side
      % walls between cells that no door connects
      \draw[wall segment] (2,0)  -- (2,-1);
      \draw[wall segment] (1,-1) -- (1,-2);
      \draw[wall segment] (3,-1) -- (3,-2);
      \draw[wall segment] (2,-2) -- (2,-3);
      \draw[wall segment] (2,-3) -- (2,-4);
      \draw[wall segment] (0,-3) -- (1,-3);
      \draw[wall segment] (1,-2) -- (2,-2);
      \draw[wall segment] (2,-3) -- (3,-3);
      \draw[wall segment] (3,-1) -- (4,-1);
      \node[exitorange, font=\scriptsize\bfseries] at (2.5,-1.5) {\textsf{E}};
      \node[exitorange, font=\scriptsize, right] at (4.15,-1.5) {the exit cell};
    \end{tikzpicture}
    \caption{The maze as walls. The area is enclosed; the exit is a cell, not a
    gap in the outer wall.}
    \label{fig:maze-as-walls}
  \end{subfigure}\hfill
  \begin{subfigure}[b]{0.46\textwidth}
    \centering
    \begin{tikzpicture}[scale=1.27]
      \foreach \x in {0,1,2,3} \foreach \y in {0,1,2,3}
        \node[node vertex] (v\x\y) at (\x,-\y) {};
      \node[exit vertex] at (2,-1) {};   % overdrawn: the one exit
      \foreach \a/\b in {v00/v10, v20/v30, v11/v21, v02/v12, v22/v32,
                         v03/v13, v23/v33,
                         v00/v01, v01/v02, v10/v11, v12/v13,
                         v20/v21, v21/v22, v31/v32, v32/v33}
        \draw[door] (\a) -- (\b);
      \node[exitorange, font=\scriptsize, below=1pt] at (2,-1.1) {\textsf{E}};
    \end{tikzpicture}
    \caption{The same maze as a spanning tree of the grid graph.}
    \label{fig:maze-as-graph}
  \end{subfigure}
  \caption{A maze is not a picture but a graph. Cells become vertices, doors
  become edges, and a wall is simply an edge of the grid graph that is absent.
  The exit is a vertex marked on the finished tree --- here an interior one,
  with degree $3$ --- and neither its position nor its degree affects anything
  the construction does.}
  \label{fig:maze-and-its-graph}
\end{figure}
